% Part: sets-functions-relations
% Chapter: functions
% Section: composition

\documentclass[../../../include/open-logic-section]{subfiles}

\begin{document}

\olfileid{sfr}{fun}{cmp}
\olsection{વિધેયોનું સંયોજન}

\begin{explain}
\oliflabeldef{sfr:fun:inv:sec}{આપણે \olref[inv]{sec}માં
જોયું કે !!a{bijection}~$f$નું પ્રતિવિધેય~$f^{-1}$
પોતે એક વિધેય છે. વિધેયો પરની બીજી ક્રિયા સંયોજન છે:
આપણે}{આપણે} બે વિધેયો $f$ અને~$g$નું સંયોજન કરીને,
એટલે કે પહેલાં $f$ અને પછી $g$ લાગુ પાડીને, નવું વિધેય
વ્યાખ્યાયિત કરી શકીએ. અલબત્ત, વિસ્તાર અને પ્રદેશ મેળ ખાતા
હોય તો જ આ શક્ય છે: $f$નો વિસ્તાર $g$ના પ્રદેશનો
ઉપગણ હોવો જોઈએ. \oliflabeldef{sfr:rel:ops:sec}{વિધેયો
પરની આ ક્રિયા \olref[rel][ops]{sec}માં સંબંધો પરની
સાપેક્ષ ગુણાકારની ક્રિયાને અનુરૂપ છે.}{}

સંયોજનનો ખ્યાલ સમજાવવામાં ચિત્ર મદદરૂપ થઈ શકે.
\olref{fig:composition}માં આપણે બે વિધેયો
$f \colon A \to B$ અને $g \colon B \to C$
તથા તેમનું સંયોજન~$(\comp{f}{g})$ દર્શાવીએ છીએ.
વિધેય $(\comp{f}{g}) \colon A \to C$ એ $A$ના દરેક
!!{element}ને $C$ના !!a{element} સાથે જોડે છે.
$A$નો !!a{element} એ $C$ના કયા !!{element}
સાથે જોડાય છે તે આ રીતે નક્કી કરીએ: આગત $x \in
A$ આપેલું હોય ત્યારે પહેલાં વિધેય $f$ને $x$ પર લાગુ પાડો.
તે કોઈક $f(x)
= y \in B$ આપશે; પછી $g$ને $y$ પર લાગુ પાડો.
તે કોઈક $g(f(x)) = g(y) = z \in C$ આપશે.
\begin{figure}
  \olasset[2\olphotowidth]{assets/diagrams/composition.tikz}
  \caption{બે વિધેયો $f$ અને~$g$નું સંયોજન $g \circ f$.}
  \ollabel{fig:composition}
\end{figure}
\end{explain}

\begin{defn}[સંયોજન]
ધારો કે $f\colon A \to B$ અને $g\colon B \to C$
વિધેયો છે. $f$નું $g$ સાથેનું \emph{સંયોજન}
એ $\comp{f}{g} \colon A \to C$ છે, જ્યાં
$(\comp{f}{g})(x) = g(f(x))$.
\end{defn}

\begin{ex}
વિધેયો $f(x) = x + 1$ અને $g(x) = 2x$ લો.
$(\comp{f}{g})(x) = g(f(x))$ હોવાથી, દરેક આગત~$x$
માટે પહેલાં તેનું અનુગામી લેવું પડે અને પછી પરિણામને
બે વડે ગુણવું પડે. તેથી તેમનું સંયોજન
$(\comp{f}{g})(x) = 2(x+1)$ દ્વારા અપાય છે.
\end{ex}

\begin{prob}
બતાવો કે જો $f \colon A \to B$ અને $g \colon B \to C$
બંને !!{injective} હોય, તો $\comp{f}{g}\colon A \to C$
પણ !!{injective} છે.
\end{prob}

\begin{prob}
બતાવો કે જો $f \colon A \to B$ અને $g \colon B \to C$
બંને !!{surjective} હોય, તો $\comp{f}{g}\colon A \to C$
પણ !!{surjective} છે.
\end{prob}

\begin{prob}
ધારો કે $f \colon A \to B$ અને $g \colon B \to C$ છે.
બતાવો કે $\comp{f}{g}$નો આલેખ $R_f \mid R_g$ છે.
\end{prob}

\end{document}
